\chapter{附录C：主要来源与声明}

\section{公司公告}

主要官方材料包括六家公司2025年年度报告、2026年第一季度报告和2026年半年度业绩预告，以及恩捷股份限制性股票注销、中国船舶分红除息和江波龙限制性股票归属相关公告。公司预告均为未经审计估计，不能替代半年报。

\begin{itemize}
  \item 中国船舶（600150）：2025年报、2026Q1、2026H1业绩预告及分红实施股本公告；
  \item 江波龙（301308）：2025年报、2026Q1、2026H1业绩预告及股本变动公告；
  \item 恩捷股份（002812）：2025年报、2026Q1、2026H1业绩预告及限制性股票注销公告；
  \item 盛新锂能（002240）：2025年报、2026Q1、2026H1业绩预告；
  \item 天华新能（300390）：更正后的2025年报、2026Q1、2026H1业绩预告；
  \item 雅化集团（002497）：2025年报、2026Q1、2026H1业绩预告。
\end{itemize}

\section{原始券商报告}

本报告归档并阅读14份原始PDF。正式模型的限权外部锚只引用东吴证券《雅化集团：26H1业绩预告点评》（2026年7月7日，目标42元）；东吴证券《恩捷股份：26H1业绩预告点评》（2026年7月10日，目标103元）因公司级经营证据阻断，仅作零权敏感性。其他报告用于盈利预测、经营假设与风险比较，不用于点目标。

中国船舶覆盖来自五份公开原始报告；江波龙两份；盛新锂能四份；天华新能一份。公开集合全为正面评级，存在明显选择偏差。

\section{行情与市场环境}

个股价格使用2026年7月22日收盘后结构化行情，并以腾讯行情交叉核对。市场环境引用当日公开收盘报道：上证3,867.03点、深证成指下跌1.42\%、创业板下跌3.23\%、两市成交约2.65万亿元、超过3,800只股票下跌。

\section{使用声明}

本报告是截至2026年7月22日的研究快照，不构成证券买卖、仓位管理或收益承诺。目标价与概率均依赖公开信息和明确假设。市场、公司经营、会计口径与券商预测均可能变化；任何决策前应重新获取最新公告、行情与风险承受能力信息。

\begin{disclosurebox}[最终结论复述]
没有高置信度基准翻倍标的。雅化集团是唯一条件核心候选；恩捷股份虽有显式牛市目标，但因经营证据不足只保留零权敏感性；中国船舶是质量备选；其余标的保持周期或高风险观察。牛市目标不等于预期收益。
\end{disclosurebox}
